% !TEX root = ../../main.tex
% Author requested this recall, diagrams, and Euclidean examples on September 30, 2026.
% Editorial clarification: assigning each input an output defines a function;
% injectivity additionally requires that distinct inputs have distinct outputs.
\section{Recall: injective, surjective, and bijective maps}
\label{sec:appendix:map-properties}

Let \(f:A\to B\) be a function, also called a map. The set \(A\) is its
\emph{domain}, and \(B\) is its \emph{codomain}. Every element of \(A\)
has exactly one output in \(B\). This requirement alone does not say
whether every element of \(B\) is reached or whether two different
inputs can have the same output.

\subsection{Surjectivity}

The map \(f\) is \emph{surjective}, or \emph{onto}, if every element
of the codomain is the output of at least one input. Equivalently,
\[
  \text{for every }b\in B,\text{ there exists }a\in A
  \text{ such that }f(a)=b.
\]
Different inputs may have the same output. In the middle panel of
\autoref{fig:recall:map-properties}, every codomain element receives
an arrow, and two inputs share an output.

\subsection{Injectivity}

The map \(f\) is \emph{injective}, or \emph{one-to-one}, if distinct
inputs have distinct outputs. Equivalently,
\[
  f(a_1)=f(a_2)\quad\Longrightarrow\quad a_1=a_2.
\]
\marginnote{As a mnemonic, think of the \emph{in} in \emph{injective}
  as a reminder to check the \emph{inputs}. Different inputs must
  remain distinguishable through their outputs.}
Each codomain element has at most one input mapping to it. Some
codomain elements may have no such input, as the top panel of
\autoref{fig:recall:map-properties} shows.

\subsection{Bijectivity}

The map \(f\) is \emph{bijective} if it is both injective and surjective.
Every codomain element then has exactly one input mapping to it.
This gives a one-to-one correspondence between the entire domain
and codomain, as in the bottom panel of
\autoref{fig:recall:map-properties}. Bijectivity imposes both
conditions, so it is stronger than either condition alone.

\begin{figure}[htbp]
  \centering
  % Original finite-set schematics for the author's requested map recall.
\begin{tikzpicture}[
  x=1mm,y=1mm,font=\footnotesize,line width=0.45pt,
  every node/.style={inner sep=2pt},
  map arrow/.style={line width=0.35pt,-{Latex[length=1.6mm,width=1.1mm]}}
]
  \begin{scope}
    \node at (44,28) {Injective, not surjective};
    \draw[fill=black!5] (15,0) ellipse[x radius=11mm,y radius=18mm];
    \draw[fill=black!5] (73,0) ellipse[x radius=11mm,y radius=18mm];
    \node at (15,22) {\(A\)};
    \node at (73,22) {\(B\)};
    \foreach \y/\i in {9/1,0/2,-9/3} {
      \fill (15,\y) circle[radius=1.6pt];
      \node at (9,\y) {\(a_{\i}\)};
    }
    \foreach \y/\i in {9/1,3/2,-3/3,-9/4} {
      \fill (73,\y) circle[radius=1.6pt];
      \node at (79,\y) {\(b_{\i}\)};
    }
    \draw[map arrow] (15,9) -- (73,9);
    \draw[map arrow] (15,0) -- (73,3);
    \draw[map arrow] (15,-9) -- (73,-3);
  \end{scope}
  \begin{scope}[yshift=-52mm]
    \node at (44,28) {Surjective, not injective};
    \draw[fill=black!5] (15,0) ellipse[x radius=11mm,y radius=18mm];
    \draw[fill=black!5] (73,0) ellipse[x radius=11mm,y radius=18mm];
    \node at (15,22) {\(A\)};
    \node at (73,22) {\(B\)};
    \foreach \y/\i in {9/1,3/2,-3/3,-9/4} {
      \fill (15,\y) circle[radius=1.6pt];
      \node at (9,\y) {\(a_{\i}\)};
    }
    \foreach \y/\i in {9/1,0/2,-9/3} {
      \fill (73,\y) circle[radius=1.6pt];
      \node at (79,\y) {\(b_{\i}\)};
    }
    \draw[map arrow] (15,9) -- (73,9);
    \draw[map arrow] (15,3) -- (73,0);
    \draw[map arrow] (15,-3) -- (73,0);
    \draw[map arrow] (15,-9) -- (73,-9);
  \end{scope}
  \begin{scope}[yshift=-104mm]
    \node at (44,28) {Bijective};
    \draw[fill=black!5] (15,0) ellipse[x radius=11mm,y radius=18mm];
    \draw[fill=black!5] (73,0) ellipse[x radius=11mm,y radius=18mm];
    \node at (15,22) {\(A\)};
    \node at (73,22) {\(B\)};
    \foreach \y/\i in {9/1,0/2,-9/3} {
      \fill (15,\y) circle[radius=1.6pt];
      \fill (73,\y) circle[radius=1.6pt];
      \node at (9,\y) {\(a_{\i}\)};
      \node at (79,\y) {\(b_{\i}\)};
      \draw[map arrow] (15,\y) -- (73,\y);
    }
  \end{scope}
\end{tikzpicture}

  \caption[Injective, surjective, and bijective maps.]{Three maps between finite sets. Each dot is an element.
    Top: an injective map leaves an output unused.
    Middle: a surjective map reaches every output, but two inputs share
    an output. Bottom: a bijective map reaches every output exactly once.
    These finite-set diagrams illustrate the definitions; the ellipses
    are not pictures of vector subspaces.}
  \label{fig:recall:map-properties}
\end{figure}

\subsection{Examples in Euclidean spaces}

The same definitions apply to maps between spaces such as \(\R^2\)
and \(\R^3\). The domain and codomain are part of each example.

\paragraph{Injective: including a plane in three-dimensional space.}
Consider
\[
  Z:\R^2\to\R^3,\qquad Z(x,y)=(x,y,0).
\]
Different pairs retain their distinct first two coordinates, so \(Z\)
is injective. Its image is the \(xy\)-plane. It is not surjective onto
\(\R^3\), since an output with a nonzero third coordinate is never reached.
\autoref{fig:recall:plane-inclusion} shows a pair and its image in that
plane.

\begin{figure}[htbp]
  \centering
  % Schematic inclusion of R^2 as the xy-plane in R^3.
\begin{tikzpicture}[
  x=1mm,y=1mm,font=\footnotesize,line width=0.45pt,
  every node/.style={inner sep=2pt},
  axis/.style={-{Latex[length=1.6mm,width=1.1mm]}},
  map arrow/.style={line width=0.35pt,-{Latex[length=1.6mm,width=1.1mm]}}
]
  \node at (44,39) {\(Z(x,y)=(x,y,0)\)};
  \node at (10,29) {\(\R^2\)};
  \node at (72,29) {\(\R^3\)};
  \draw[axis] (0,0) -- (24,0) node[right] {\(x\)};
  \draw[axis] (0,0) -- (0,22) node[above] {\(y\)};
  \draw[gray] (0,0) -- (12,8);
  \fill (12,8) circle[radius=1.6pt];
  \begin{scope}[shift={(52,0)},x={(8mm,-3mm)},y={(8mm,3mm)},z={(0mm,9mm)}]
    \filldraw[fill=black!10] (0,0,0) -- (2,0,0) -- (2,2,0)
      -- (0,2,0) -- cycle;
    \draw[axis] (0,0,0) -- (2.7,0,0) node[below] {\(x\)};
    \draw[axis] (0,0,0) -- (0,2.7,0) node[above right] {\(y\)};
    \draw[axis] (0,0,0) -- (0,0,2.4) node[above] {\(z\)};
    \draw[gray] (0,0,0) -- (1.5,1,0);
    \fill (1.5,1,0) circle[radius=1.6pt];
  \end{scope}
  \draw[map arrow] (12,8) .. controls (30,19) and (49,13) .. (72,-1.5);
\end{tikzpicture}

  \caption[Plane inclusion into three-dimensional space.]{Schematic inclusion \(Z:\R^2\to\R^3\). The shaded region
    depicts part of the \(xy\)-plane, which continues beyond the drawing.
    The marked input maps to a point in this plane. Distinct input
    coordinates remain distinct after inclusion.}
  \label{fig:recall:plane-inclusion}
\end{figure}

\paragraph{Surjective: forgetting one coordinate.}
Consider
\[
  U:\R^3\to\R^2,\qquad U(x,y,z)=(y,z).
\]
Every pair \((y,z)\) is attained, for instance from \((0,y,z)\), so
\(U\) is surjective. It is not injective, since changing \(x\) leaves
the output unchanged. \autoref{fig:recall:coordinate-projection}
shows inputs differing in the \(x\)-direction sharing an output.

\begin{figure}[htbp]
  \centering
  % Schematic coordinate projection R^3 -> R^2, (x,y,z) -> (y,z).
\begin{tikzpicture}[
  x=1mm,y=1mm,font=\footnotesize,line width=0.45pt,
  every node/.style={inner sep=2pt},
  axis/.style={-{Latex[length=1.6mm,width=1.1mm]}},
  map arrow/.style={line width=0.35pt,-{Latex[length=1.6mm,width=1.1mm]}}
]
  \node at (44,39) {\(U(x,y,z)=(y,z)\)};
  \node at (10,29) {\(\R^3\)};
  \node at (68,29) {\(\R^2\)};
  \begin{scope}[x={(8mm,-3mm)},y={(8mm,3mm)},z={(0mm,9mm)}]
    \draw[axis] (0,0,0) -- (3,0,0) node[below] {\(x\)};
    \draw[axis] (0,0,0) -- (0,2.6,0) node[below right] {\(y\)};
    \draw[axis] (0,0,0) -- (0,0,2.4) node[above] {\(z\)};
    \draw[gray,dashed] (0,0.8,1.8) -- (2.7,0.8,1.8);
    \fill (1,0.8,1.8) circle[radius=1.6pt];
    \fill (2,0.8,1.8) circle[radius=1.6pt];
  \end{scope}
  \draw[axis] (56,0) -- (80,0) node[right] {\(y\)};
  \draw[axis] (56,0) -- (56,22) node[above] {\(z\)};
  \fill (64,18) circle[radius=1.6pt];
  \draw[gray] (56,0) -- (64,18);
  \draw[map arrow] (14.4,13.8)
    .. controls (32,22) and (46,22) .. (64,18);
  \draw[map arrow] (22.4,10.8)
    .. controls (38,13) and (48,15) .. (64,18);
\end{tikzpicture}

  \caption[Projection onto two coordinates.]{Schematic projection \(U:\R^3\to\R^2\). The marked inputs
    lie on a line parallel to the \(x\)-axis and have the same \(y,z\)
    coordinates. They map to the same pair. Every pair in the codomain
    is attained.}
  \label{fig:recall:coordinate-projection}
\end{figure}

\paragraph{Bijective: swapping two coordinates.}
Consider
\[
  B:\R^2\to\R^2,\qquad B(x,y)=(y,x).
\]
Each output \((a,b)\) has exactly one preimage, namely \((b,a)\).
Thus \(B\) is bijective. Swapping coordinates again recovers the
original input. \autoref{fig:recall:coordinate-swap} shows the
reflection across the line \(x=y\).

\begin{figure}[htbp]
  \centering
  % Schematic coordinate swap R^2 -> R^2, (x,y) -> (y,x).
\begin{tikzpicture}[
  x=1mm,y=1mm,font=\footnotesize,line width=0.45pt,
  every node/.style={inner sep=2pt},
  axis/.style={-{Latex[length=1.6mm,width=1.1mm]}},
  map arrow/.style={line width=0.35pt,-{Latex[length=1.6mm,width=1.1mm]}}
]
  \node at (44,39) {\(B(x,y)=(y,x)\)};
  \node at (10,29) {\(\R^2\)};
  \node at (68,29) {\(\R^2\)};
  \foreach \origin in {0,56} {
    \begin{scope}[xshift=\origin mm]
      \draw[axis] (0,0) -- (24,0) node[right] {\(x\)};
      \draw[axis] (0,0) -- (0,22) node[above] {\(y\)};
      \draw[gray,dashed] (0,0) -- (20,20);
    \end{scope}
  }
  \draw[gray] (0,0) -- (16,8);
  \draw[gray] (56,0) -- (64,16);
  \fill (16,8) circle[radius=1.6pt];
  \fill (64,16) circle[radius=1.6pt];
  \draw[map arrow] (16,8) .. controls (32,18) and (48,18) .. (64,16);
\end{tikzpicture}

  \caption[Swapping two coordinates.]{Schematic coordinate swap in \(\R^2\). The dashed diagonal
    in each coordinate system is \(x=y\). The marked point and its
    image have exchanged coordinates. Each output has exactly one input.}
  \label{fig:recall:coordinate-swap}
\end{figure}
